\section{Keyboard and Preferences}

\subsection{Keys}

\begin{req}{KEY-1}
The interface can be operated entirely from the keyboard.  Navigation keys (arrows,
Home, End, Page Up, Page Down, Space) are matched by the physical key, so they work
whatever modifier is held; so are the Alt shortcuts below, because macOS changes the
character that Option and a letter produce (Option+C types ``\copyright'').  Other
letters and punctuation are matched by the character they produce, so they follow the
user's keyboard layout.  No shortcut acts while
the user is typing in a text box, except Escape, which leaves it.
\end{req}

\begin{req}{KEY-2}
The following keys are defined.

\begin{center}
\begin{tabular}{@{}p{0.42\linewidth}p{0.5\linewidth}@{}}
\toprule
Key & Action \\
\midrule
\code{Up}, \code{Down}, \code{PgUp}, \code{PgDn}, \code{Home}, \code{End} & Move the selection \\
\code{Enter} or \code{Right} & Open the selected folder, selecting its first entry (a search result opens in its folder, with it selected) \\
\code{Left} or \code{Backspace} & Up one folder, reselecting the one just left \\
\code{Space} & Show or hide the preview pane \\
\code{/} & Filter this folder \\
A letter or digit & Jump to the next entry whose name starts with what has been typed (\rid{KEY-3}) \\
\code{Alt+S} & Search subfolders by name (Enter runs, Escape clears) \\
\code{Alt+G} & Go to a path \\
\code{Alt+C} & Copy the selected path (or the folder's) \\
\code{Escape} & Clear the search, close a menu, leave a text box \\
\code{Alt+Left}, \code{Alt+Right} on a column header & Reorder the column \\
\bottomrule
\end{tabular}
\end{center}
\end{req}

\begin{req}{KEY-3}
Type-ahead: typing a letter or digit selects the next entry, after the current one
and wrapping around, whose name starts with it, compared as the filter compares names
(\rid{LST-8}).  Letters typed within 800~ms of each other extend the prefix; repeating
the same letter cycles through the entries that start with it.  No plain letter is
reserved for a shortcut, so every name can be reached this way.
\end{req}

\subsection{Mouse}

\begin{req}{MOUSE-1}
A plain left click on a row selects it, exactly as moving to it with the keyboard
would; a click on a folder also opens it.  A plain click on a file never downloads it.
A modified click (with Command, Control, Shift or Alt) or the browser's context menu on
a file name still offers the browser's own download, which is the explicit way to
download (\rid{API-4}).
\end{req}

\subsection{Preferences}
\label{sec:prefs}

\begin{req}{PREF-1}
The following are the user's preferences and are kept in the browser's local storage,
per browser and never on the server or on disk elsewhere: whether the preview pane is
open and how wide it is, whether hover peek is on, ``Folders first'', ``Match case'',
and the layout of the columns (order, widths, which are hidden).
\end{req}

\begin{req}{PREF-2}
Preferences are optional.  If local storage is missing, blocked, corrupt or from an
older version, each value falls back to its default on its own and the interface works
normally.
\end{req}
